
\documentclass{article}
\usepackage{colortbl}
\usepackage{makecell}
\usepackage{multirow}
\usepackage{supertabular}

\begin{document}

\newcounter{utterance}

\centering \large Interaction Transcript for game `matchit\_ascii', experiment `different\_grid', episode 2 with selfrefine{-}Qwen3.5{-}9B.
\vspace{24pt}

{ \footnotesize  \setcounter{utterance}{1}
\setlength{\tabcolsep}{0pt}
\begin{supertabular}{c@{$\;$}|p{.15\linewidth}@{}p{.15\linewidth}p{.15\linewidth}p{.15\linewidth}p{.15\linewidth}p{.15\linewidth}}
    \# & $\;$A & \multicolumn{4}{c}{Game Master} & $\;\:$B\\
    \hline

    \theutterance \stepcounter{utterance}  
    & & \multicolumn{4}{p{0.6\linewidth}}{
        \cellcolor[rgb]{0.9,0.9,0.9}{
            \makecell[{{p{\linewidth}}}]{
                \texttt{\tiny{[P1$\langle$GM]}}
                \texttt{You are participating in a collaborative guessing game.} \\
\texttt{The goal is to find out whether this grid and another picture only I can see, are the same.} \\
\texttt{Please describe your grid first. Then, I will provide my description and we can ask each other questions about the grids to figure out whether they are the same. Now start your short grid description with "DESCRIPTION:" followed by the description. Do not add anything else. This is your grid:} \\
\texttt{X ▢ ▢ ▢ X} \\
\texttt{▢ X ▢ ▢ X} \\
\texttt{▢ ▢ X X X} \\
\texttt{▢ X ▢ ▢ X} \\
\texttt{X ▢ ▢ ▢ X} \\
            }
        }
    }
    & \\ \\

    \theutterance \stepcounter{utterance}  
    & \multicolumn{4}{p{0.6\linewidth}}{
        \cellcolor[rgb]{0.95,0.95,0.95}{
            \makecell[{{p{\linewidth}}}]{
                \texttt{\tiny{[P2$\rightarrow$P1]}}
                \texttt{{[}DRAFT{]}} \\
\texttt{DESCRIPTION: A 5x5 grid where the border cells (first and last rows, first and last columns) are all 'X'. The interior 3x3 area contains a pattern of 'X's and empty squares '▢', specifically with 'X's at positions (2,2), (3,3), (4,4) relative to the top{-}left of the interior, forming a diagonal line from top{-}left to bottom{-}right within the center, while the rest of the interior is '▢'.} \\
            }
        }
    }
    & & \\ \\

    \theutterance \stepcounter{utterance}  
    & \multicolumn{4}{p{0.6\linewidth}}{
        \cellcolor[rgb]{0.95,0.95,0.95}{
            \makecell[{{p{\linewidth}}}]{
                \texttt{\tiny{[P2$\rightarrow$P1]}}
                \texttt{{[}SELF{-}FEEDBACK{]}} \\
\texttt{FORMAT: ok} \\
\texttt{RULES: ok} \\
\texttt{STRATEGY: problem} \\
\texttt{REASON: The description incorrectly identifies the interior pattern as a diagonal line, whereas the actual grid shows 'X's at (2,2), (3,3), and (4,4) which are not on a single diagonal line relative to the 3x3 interior (the interior indices are 1{-}3, so the X's are at 1,1; 2,2; 3,3, which is a diagonal, but the description claims the rest of the interior is empty while the grid has X's at (1,3) and (2,1) relative to the full grid which correspond to interior positions (1,3) and (2,1) {-} wait, let's re{-}verify the grid).} \\
\\ 
\texttt{Re{-}evaluating the grid:} \\
\texttt{Row 1: X ▢ ▢ ▢ X} \\
\texttt{Row 2: ▢ X ▢ ▢ X} \\
\texttt{Row 3: ▢ ▢ X X X} \\
\texttt{Row 4: ▢ X ▢ ▢ X} \\
\texttt{Row 5: X ▢ ▢ ▢ X} \\
\\ 
\texttt{Interior (rows 2{-}4, cols 2{-}4):} \\
\texttt{Row 2 interior: X ▢ ▢} \\
\texttt{Row 3 interior: ▢ X X} \\
\texttt{Row 4 interior: X ▢ ▢} \\
\\ 
\texttt{The user's draft says: "X's at positions (2,2), (3,3), (4,4) relative to the top{-}left of the interior".} \\
\texttt{If interior is 3x3:} \\
\texttt{(1,1) is X (from grid row 2 col 2)} \\
\texttt{(1,2) is ▢} \\
\texttt{(1,3) is ▢} \\
\texttt{(2,1) is ▢} \\
\texttt{(2,2) is X (from grid row 3 col 3)} \\
\texttt{(2,3) is X (from grid row 3 col 4)} \\
\texttt{(3,1) is X (from grid row 4 col 2)} \\
\texttt{(3,2) is ▢} \\
\texttt{(3,3) is ▢} \\
\\ 
\texttt{The draft claims X's at (2,2), (3,3), (4,4) relative to interior. This implies a 4x4 interior or wrong indexing. Even if we assume 1{-}based indexing for the 3x3 interior:} \\
\texttt{Draft says X at (2,2), (3,3), (4,4). (4,4) doesn't exist in a 3x3.} \\
\texttt{The actual X's in the interior are at (1,1), (2,2), (2,3), (3,1).} \\
\texttt{The draft describes a diagonal line (2,2), (3,3), (4,4) which is factually incorrect for the provided grid. This} \\
            }
        }
    }
    & & \\ \\

    \theutterance \stepcounter{utterance}  
    & \multicolumn{4}{p{0.6\linewidth}}{
        \cellcolor[rgb]{0.9,0.9,0.9}{
            \makecell[{{p{\linewidth}}}]{
                \texttt{\tiny{[P1$\rangle$GM]}}
                \texttt{DESCRIPTION: A 5x5 grid where the border cells (first and last rows, first and last columns) are all 'X'. The interior 3x3 area contains 'X's at positions (1,1), (2,2), (2,3), and (3,1) relative to the top{-}left of the interior, while the remaining interior cells are '▢'.} \\
            }
        }
    }
    & & \\ \\

    \theutterance \stepcounter{utterance}  
    & & & \multicolumn{2}{p{0.3\linewidth}}{
        \cellcolor[rgb]{0.9,0.9,0.9}{
            \makecell[{{p{\linewidth}}}]{
                \texttt{\tiny{[GM$|$GM]}}
                \texttt{continue} \\
            }
        }
    }
    & & \\ \\

    \theutterance \stepcounter{utterance}  
    & & \multicolumn{4}{p{0.6\linewidth}}{
        \cellcolor[rgb]{0.9,0.9,0.9}{
            \makecell[{{p{\linewidth}}}]{
                \texttt{\tiny{[GM$\rangle$P2]}}
                \texttt{You are participating in a collaborative guessing game.} \\
\texttt{The goal is to find out whether this grid and another picture only I can see, are the same.} \\
\texttt{Please describe your grid first. Then, I will provide my description and we can ask each other questions about the grids to figure out whether they are the same. Now start your short grid description with "DESCRIPTION:" followed by the description. Do not add anything else. This is your grid:} \\
\texttt{X X X X ▢} \\
\texttt{X X X ▢ X} \\
\texttt{X X ▢ X X} \\
\texttt{X ▢ X X X} \\
\texttt{▢ X X X X} \\
            }
        }
    }
    & \\ \\

    \theutterance \stepcounter{utterance}  
    & & & \multicolumn{4}{p{0.6\linewidth}}{
        \cellcolor[rgb]{0.95,0.95,0.95}{
            \makecell[{{p{\linewidth}}}]{
                \texttt{\tiny{[P1$\rightarrow$P2]}}
                \texttt{{[}DRAFT{]}} \\
\texttt{DESCRIPTION: A 5x5 grid where the first four rows each contain exactly one empty square (▢) in a different column, and the fifth row has the empty square in the first column. Specifically, the empty squares are at positions (1,4), (2,4), (3,3), (4,2), and (5,1).} \\
            }
        }
    }
     \\ \\

    \theutterance \stepcounter{utterance}  
    & & & \multicolumn{4}{p{0.6\linewidth}}{
        \cellcolor[rgb]{0.95,0.95,0.95}{
            \makecell[{{p{\linewidth}}}]{
                \texttt{\tiny{[P1$\rightarrow$P2]}}
                \texttt{{[}SELF{-}FEEDBACK{]}} \\
\texttt{FORMAT: ok} \\
\texttt{RULES: problem} \\
\texttt{STRATEGY: problem} \\
\texttt{REASON: The answer violates the game's instruction to describe the grid first without adding extra text, and the strategy is flawed because the provided coordinates (1,4) and (2,4) are incorrect based on the visual grid where the empty squares in the first two rows are in the fourth and fifth columns respectively.} \\
            }
        }
    }
     \\ \\

    \theutterance \stepcounter{utterance}  
    & & & \multicolumn{4}{p{0.6\linewidth}}{
        \cellcolor[rgb]{0.9,0.9,0.9}{
            \makecell[{{p{\linewidth}}}]{
                \texttt{\tiny{[GM$\langle$P2]}}
                \texttt{ok} \\
            }
        }
    }
     \\ \\

    \theutterance \stepcounter{utterance}  
    & & & \multicolumn{2}{p{0.3\linewidth}}{
        \cellcolor[rgb]{0.9,0.9,0.9}{
            \makecell[{{p{\linewidth}}}]{
                \texttt{\tiny{[GM$|$GM]}}
                \texttt{abort, first word: ok, but expected DESCRIPTION:} \\
            }
        }
    }
    & & \\ \\

\end{supertabular}
}

\end{document}
